\section{序論}

同じ関数斜体を持つ非可換射影空間を、射影導来圏の段階でどこまで区別できるか。
さらに、その半直交成分だけを取り出した後にも、量子化の違いは残るか。
本稿は重み $(1,1,1,m)$ の明示的な非線形CGG量子化族について、
固定非零点の変形環、Serre関手、およびGerstenhaber作用からこの問題を調べる。
四項成分では、代数自身の滑らかな変形基底の中に、
対称性を保持する条件によって非被約な二重構造が現れることを示す。
さらに引戻し型のOre量子化を加え、非完全交叉の変形環と、全変形の明示的なAS正則実現を得る。

対象は整数 $j,r\ge0$、$m=2j+r\ge3$ と
$\alpha\in\mathbb C$、$s\in\mathbb C^*$ による
$\mathcal S_{j,r}(\alpha,s)$ である。積は次節で完全に定義する。
各代数はAS正則四次元であり、標準 $(m+3)$ 項のfull strong exceptional collectionを持つ。
全ての族の次数零関数斜体は、$\mathbb C(C)$ 上の第一Weyl斜体に同型となる。
その一方、射影圏のアーベル圏同値と有界導来圏同値は共に
\[
 \boxed{(m,j)=(m',j')}
\]
によって特徴付けられる。固定した $m$ に対し、相異なる圏は
$\lfloor m/2\rfloor+1$ 個である。分類対象は、ここで明示したCGG族である。

\subsection{全共鳴の変形環と端点の相違}

標準全列の有限自己準同型代数を $A_{j,r}$ とすると、固定非零量子化のminiversal環は
\[
 R_{A_{j,r}}\simeq
 \begin{cases}
 \mathbb C[[\gamma]]/(\gamma^{j+1})&r>0,\\
 \mathbb C[[u_1,\ldots,u_{m+2}]]&r=0
 \end{cases}
\]
となる。第一の場合は任意に高い次数で初めて現れる障害を与え、
$j\ge1$ では一次変形があるにもかかわらず、全ての一パラメータ形式変形が自明である。
第二の場合は $m=2j$ の偶数重みの端点であり、$m+2$ 次元の滑らかな変形基底が現れる。
端点の証明では、混合Poisson写像の核をHamiltonianベクトル場で尽くし、
調和多項式の分解によって余核がただ一方向であることを示す。
その方向を非零混合点の横断パラメータが消すため、全局所環が正則になる。

Poissonスキームから固定量子代数への移送には、可換基点の形式同型を
$h$ 進に小さい量子化弧へ適用し、$\mathbb C((h))$ 上で再中心化する方法を用いる。
全局所環を使うため、接空間の一致だけに依存しない。
得られるのは有限Beilinson代数、従って射影導来圏の変形環である。
無限次元の次数付き環の全ての非次数付き変形を分類しているわけではない。


\subsection{全変形の座標環実現と非完全交叉hull}

偶数端点の $m+2$ 次元の形式変形は、三変数AS正則代数 $B_{\gamma,s}$ と
混合項 $H$ から作る $B_{\gamma,s}[w;\sigma,x^{-1}[\,-,H]]$ によって全て実現できる。
正規元 $x$ の正規化自己同型が $B$ のNakayama自己同型と一致するため、
中心元 $J=xw+H$ の固定性から拡大のNakayama自己同型が恒等写像となる。

これは各複素パラメータでAS正則となる有限パラメータ族である。
その次数零関数斜体を全パラメータで分類する。
混合項 $H$ の $m+1$ 方向は関数斜体から消えるが、非冪根の場合には残る $\gamma$ を完全に区別できる。
冪根の場合には中心変数を混ぜる整数変換により、同じ位数のパラメータが同じ斜体を与える。
この族のPI集合は構成可能でなく、任意に高いPI次数のファイバーが非PI端点へ集積する。

一方、共点三直線に基づく微分Ore量子化では、固定非零点のhullが
\[
 \mathbb C[[t_1,t_2,a,b]]/(a^{m+2},ab,b^2)
\]
となる。二つの滑らかな方向はOre係数 $A,B$ の変化で実現される。
横断Artin環の長さは $m+3$、socle次元は二であり、全hullはCohen--Macaulay型二の非完全交叉環である。
従って本稿は、一変数の高次障害、滑らかな全実現、および二変数の障害の相互作用を、
実際の非可換重み付き射影圏で比較できる。

\subsection{Serre関手と内在的な反標準環}

通常のCalabi--Yau代表に対して、反標準Hochschild--Serre環の全体を
$Z(\mathcal S)^{(m+3)}$ と同定する。
$r\ge2,r=1,r=0$ に応じて、この環は定数環、$D_{m+4}$ 型、$A_{m+2}$ 型となる。
正次数部分がある場合のstacky Projは、それぞれ
$(2,2,m+2)$ 型のorbifold曲線と $\mathbb P_{\mathrm{stack}}(2,m+1)$ である。
中心から自然変換を作るだけでなく、正twist双加群の全ての準同型が中心から来ることを証明する。
この全体性によって、中心の表示を射影圏に内在する不変量として用いられる。

中心の部分解消、central reduction、primitive idealについては、
独立に読める公表稿XI\cite{VKPrimitive36}にまとめる。
本稿では導来分類に必要な中心と反標準環の計算を扱う。

\subsection{同じ滑らかな変形基底を持つ二つの量子VK圏}

重み三では $(j,r)=(0,3),(1,1)$ の二つの非零族がある。
標準四項 $(0,1,3,4)$ を切り出した五十六次元代数を $B_0,B_1$ と書く。
両者は
\[
 \dim HH^\bullet(B_i)=(1,1,35,39),\qquad
 R_{B_i}\simeq\mathbb C[[z_1,\ldots,z_{35}]]
\]
を満たす。従って、Hochschild群の次元も、miniversal基底の抽象的な形式環も一致する。
それでも、二つの導来圏は同値ではない。

区別に使うのは、一次元空間 $HH^1(B_i)$ の非零元 $\xi$ による
\[
 [\xi,-]:HH^2(B_i)\longrightarrow HH^2(B_i)
\]
の冪零Jordan型である。二つの場合のブロックはそれぞれ
\[
 (13,10,7,5),\qquad (11,9,6,5,3,1)
\]
であり、第一冪の階数だけでも $31$ と $29$ に分かれる。
生成元を非零スカラー倍してもこのJordan型は変わらず、
Gerstenhaber括弧の導来不変性から、これは標準例外列の選択によらない区別となる。
可換基点も含めると、重み三の本CGG族の四項成分は正確に三つの導来同値類に分かれる。

標準六項補完の変形についても、二つの四項成分の滑らかな基底の中で、
補完像はそれぞれ被約点と長さ二の非被約点となる。
この補完像の相違と、Gerstenhaber作用による内在的な非同値証明は別の主張である。
任意の別の例外列による補完可能性まで分類したものではない。


さらに、相対Hochschild複体を縮約すると、対称性の消失を一列の微分で記述できる。
一般変形のHochschild次元は $(1,0,34,39)$ となる。
外的導分を保持する部分の接空間は四次元と六次元であり、
二次近傍はそれぞれ無関係な四変数と、六変数に一つの平方関係を課したものとなる。
代数自身の変形は無障害でも、基点の外的導分を保持する条件には二次障害が現れる。
原始射に支持を持つ積補正を用いると、二次障害の零方向は全次数で同時に積分できる。
さらに平方零方向を含む全次数の族を構成することで、対称性保持部分の全環は
\[
 R_{Z_{B_0}}\simeq\mathbb C[[u_1,\ldots,u_4]],\qquad
 R_{Z_{B_1}}\simeq
       \mathbb C[[u_1,\ldots,u_5,\gamma]]/(\gamma^2)
\]
と完全に決まる。
第二の環は滑らかな五次元部分の二重構造であり、
基底全体で二十四乗が零となる導分、従って $\mathbf G_a$ 作用を持つ。
その構成では、有理係数の有限行列から得た解について、
全ての結合則と導分方程式を多項式恒等式として直接確認する。
有限jetから全完備局所環を推測しているのではない。

この二重構造には、純粋な三変数部分からの説明もある。
全ての変形と基底変換を $w$ フィルトレーションを保つ形へ還元し、
純粋部分への制限の相対変形空間が三十二次元で滑らかであることを示す。
純粋部分自体は、固定した四対象の合成の像から復元できる。
その変形基底は滑らかな三次元、対称性保持部分は滑らかな曲線
$\operatorname{Spf}\mathbb C[[\gamma]]$ である。
原始射を接合すると、この曲線は $B_0$ では原点に、
$B_1$ では原点の平方零近傍に制限される。
従って上の第二の同型は、純粋曲線の環 $\mathbb C[[\gamma]]$ 上で取ることができる。

通常の三十五変数hullの中でも形式座標を選び、対称性保持イデアルを
\[
 I_{B_1}=(z_1,\ldots,z_{29},\gamma^2)
\]
と書ける。標準六項への補完像は、この二重構造の横断長さ二の切片となる。
その接線を初速度とするどの積分弧でも外的導分は二次で失われ、
相対Hochschild群には長さ二の捩れが現れる。
さらに一般の弧の捩れの長さも、$z_i$ の付値と $\gamma$ の付値の二倍の最小値で決まる。

\subsection{全重みの通常hullの分解と障害の自然な比較}

四項成分の純粋三変数部分が残る理由は、重み三に固有ではない。
全ての $m=2j+r\ge3$、$j,r\ge0$ で、任意のArtin変形における純粋部分は
合成の像から一意かつ自然に復元できる。
純粋部分の変形とその同定を固定した相対変形は、
\[
 k_m=\frac{m^2+9m+28}{2}
\]
次元の滑らかな形式空間となり、相対自己同型を持たない。
証明は任意の間隔に有効な反対称化と、CGG積の有限フィルトレーションによる。
さらに、三本の隣接射を反転する普遍局所化と、半古典Poisson構造の安定化群から
$HH^1(N_m)=\mathbb C[\log(1+N)]$ を全ての重みで証明する。
この生成導分は各四項成分へ持ち上がるため、通常miniversal環にも
\[
 R_{B_{j,r}}\simeq R_{N_m}[[u_1,\ldots,u_{k_m}]]
\]
という分解が成立する。特に、同じ重みの全共鳴は偶数端点を含めて同型な通常hullを持つ。
この同型は導分の作用や対称性保持部分を同一視しない。
一般の $m$ に対する純粋部分の環 $R_{N_m}$ 自体は、まだ決定していない。

また、六項側の結合性障害と四項側の対称性障害の対応を、cochainの段階で説明する。
積 $\mu$ と導分 $\delta$ の同時変形を記録する複体を
\[
 \mathcal P^n=C^n\oplus C^{n-1},\qquad
 \partial(f,g)=(df,[\delta,f]-dg)
\]
と置く。これはAssDer対の既存の変形理論\cite{DasMandal2026}の複体である。
全充満部分圏への制限はこの複体と両立し、
外側の結合性障害が内側で境界になったとき、
その増強された障害は内側の $\operatorname{coker}[\delta,-]$ の一意の類へ移る。
標準六項補完では、これにより一次元の結合性障害線から
内在的対称性障害への自然な非零線形写像を得る。
この比較は通常の $HH^3$ 全体を対象とする写像ではなく、
導分の補正を記録した障害と、その標準補完の障害線上の帰結である。

\subsection{先行理論と証明の構成}

CGG積と可解Lie作用による量子化にはLecoutre--Sierra\cite{LecoutreSierra2019}、
AS正則代数の標準列とSerre双対性にはMinamoto--Mori\cite{MinamotoMori2011}、
次数付きtwistにはZhang\cite{Zhang1996}の理論を用いる。
Hochschild複体とそのGerstenhaber構造の導来不変性はKeller\cite{KellerDerived2003}、
dg Lie代数の変形関手の不変性はManetti\cite{Manetti2005}による。
Hochschild--Serre環の一般構成はLin--Zhang\cite{LinZhangSerre2025}に帰する。

本稿の新しい部分は、この全重み族の固定非零点の変形環、
反標準対角の全体性と具体的表示、それらによる族内同値分類、
ならびに量子VK四項成分の変形とGerstenhaber作用、
全重みの純粋部分への還元と通常hullの分解、対称性保持部分の全非被約構造の決定、
および導分付き障害を使った標準補完の自然な比較である。
可換基点のhullと $r>0$ のPoisson局所環は、公表稿IX\cite{VKPoisson34}を入力とする。
偶数端点の全局所環、全変形のOre実現、引戻し型の固定hullも本稿で証明する。
Ore拡大のNakayama公式はLiu--Wang--Wu\cite{LWWEndpoint37}、中心固定性はReyes--Rogalski\cite{ReyesRogalski2022}による。
四項成分のJordan型は、有理数上の有限Hochschild複体から厳密に決定し、
商座標を使う方法と全コサイクル空間を使う方法で独立に検証する。

公表稿I--IXのLO型関係代数、Prym幾何、相対potentialの問題は、
本稿の明示CGG族の分類とはそれぞれ別の問題として保持する。
